\section{Methodology}
\label{sec:methodology}

This section details the full modeling pipeline: target construction, a three-stage causal preprocessing pipeline, two alternative feature engineering strategies, model specifications, a walk-forward backtesting protocol, and evaluation procedures. Figure~\ref{fig:pipeline} provides a schematic overview.\footnote{All transformations are applied in a strictly causal (non-anticipative) manner: every quantity computed at time $t$ uses only information available at or before $t-1$.}

\subsection{Realized Volatility}
\label{sec:rv_definition}

The forecasting target is realized variance (RV) computed at the 30-minute bar level from 5-minute intraday log-returns. For the $j$-th 30-minute bar on trading day $d$, realized volatility is defined as
\begin{equation}
    RV_{d,j} = \sum_{i=1}^{M} r_{d,j,i}^2,
    \label{eq:rv}
\end{equation}
where $r_{d,j,i}$ denotes the $i$-th 5-minute log-return within the bar and $M = 6$ is the number of non-overlapping 5-minute intervals per 30-minute bar. Because data are recorded over 24-hour trading days (including pre-market, regular, and after-hours sessions), each day contains $48$ bars of 30 minutes each. For notational convenience, we index observations sequentially as $RV_t$ where $t$ enumerates consecutive 30-minute bars across all trading days. The forecasting objective is to predict the one-step-ahead realized volatility $RV_{t+1}$.

\subsection{Data Preprocessing}
\label{sec:preprocessing}

Raw realized volatility and exogenous variables undergo a three-stage transformation pipeline applied in the following fixed order:
\begin{equation}
    \text{Raw series} \;\xrightarrow{\text{Stage 1}}\; \text{Diurnal adjustment} \;\xrightarrow{\text{Stage 2}}\; \text{Semantic transform} \;\xrightarrow{\text{Stage 3}}\; \text{Rolling winsorization}.
    \label{eq:pipeline}
\end{equation}
The pipeline is designed to (i) remove predictable intraday patterns, (ii) stabilize distributional properties, and (iii) limit the influence of extreme observations. Importantly, at evaluation time only Stages 1 and 2 are inverted to recover raw-scale forecasts; Stage 3 (winsorization) is a lossy operation that is not reversed.

\subsubsection{Stage 1: Diurnal Adjustment}
\label{sec:diurnal}

Intraday volatility exhibits a well-documented U-shaped pattern, with elevated levels near the market open and close. To remove this deterministic seasonality, we divide each variable by a rolling time-of-day baseline. Let $s(t) \in \{1, \ldots, 48\}$ denote the time-of-day slot for observation $t$. For non-negative variables such as $RV_t$, the diurnal baseline is
\begin{equation}
    \bar{RV}_{t \mid s} = \frac{1}{|\mathcal{W}_t|} \sum_{j \in \mathcal{W}_t} RV_j,
    \label{eq:diurnal_mean}
\end{equation}
where $\mathcal{W}_t = \{j : j < t,\; s(j) = s(t)\}$ denotes the set of up to $W = 20$ most recent observations sharing the same time slot, with a minimum of 5 observations required before the baseline becomes valid. The one-period shift ($j < t$, not $j \leq t$) ensures strict causality. The adjusted series is then
\begin{equation}
    \widetilde{RV}_t = \frac{RV_t}{\bar{RV}_{t \mid s}}.
    \label{eq:diurnal_adj}
\end{equation}
For signed variables (e.g., autocovariance), the baseline is instead the rolling standard deviation within each slot, and the adjustment divides by this rolling standard deviation. Observations in the initial warm-up period (before 5 same-slot observations have accumulated) have their baseline set to 1, effectively passing through unadjusted; these observations are subsequently consumed by the training buffer and do not enter the test set.

\subsubsection{Stage 2: Semantic Transforms}
\label{sec:semantic_transforms}

Following diurnal adjustment, we apply variance-stabilizing transforms chosen according to the distributional properties of each variable class:
\begin{itemize}
    \item Squared-return and volatility measures ($RV$, bipower variation, turnover, effective spread): $x \mapsto \sqrt{x}$.
    \item Autocovariance: $x \mapsto \mathrm{sign}(x)\sqrt{|x|}$.
    \item Cubic returns: $x \mapsto x^{1/3}$.
    \item Quartic returns: $x \mapsto x^{1/4}$.
    \item Implied volatility measures (VIX, VVIX, VIX3M): $x \mapsto \log(x)$.
\end{itemize}
These transforms compress the right tail and reduce heteroskedasticity, which is particularly beneficial for linear models. The square-root transform applied to $RV$ is especially important: since models are trained on the $\sqrt{RV}$ scale, a bias correction (Duan's smearing estimator, Section~\ref{sec:duan}) is required when converting forecasts back to the raw variance scale.

\subsubsection{Stage 3: Rolling Winsorization}
\label{sec:winsorization}

After transformation, each series is clipped to rolling 5th and 95th percentile bounds computed over a trailing window of $W_{\mathrm{clip}} = 240$ observations (approximately 5 trading days):
\begin{equation}
    x_t^{w} = \mathrm{clip}\!\left(x_t;\; q_{0.05}^{(t)},\; q_{0.95}^{(t)}\right),
    \label{eq:winsor}
\end{equation}
where $q_{\alpha}^{(t)}$ denotes the $\alpha$-quantile of the 240-observation trailing window ending at $t-1$. The one-period shift ensures that quantile bounds are computed from past data only. Because winsorization is a lossy compression, it cannot be inverted at evaluation time; forecasts are mapped back to the raw scale using only the inverse of Stages 1--2 (Section~\ref{sec:duan}).

\subsection{Feature Engineering}
\label{sec:features}

We consider two alternative strategies for constructing lagged features from the preprocessed series: a hand-crafted HAR lag structure and a data-driven PCA compression. Both strategies transform each variable's history into a fixed-dimensional feature vector suitable for regression. A third, nonlinear construction --- a Laplacian eigenmap of delay-embedded \emph{residual} trajectories --- is deferred to Section~\ref{sec:spectral_knn}, since it is defined relative to a fitted baseline rather than on the preprocessed series directly.

\subsubsection{HAR Features}
\label{sec:har_features}

The Heterogeneous Autoregressive (HAR) model of \citet{Corsi2009} captures the multi-scale persistence of volatility through aggregated lagged components. We adopt a geometric base-5 lag structure adapted to the 30-minute sampling frequency, yielding the lag set
\begin{equation}
    \mathcal{L} = \{1, 5, 25, 125, 625, 3125\},
    \label{eq:har_lags}
\end{equation}
which corresponds approximately to horizons of 30 minutes, 2.5 hours, half a trading day, 2.6 trading days, 13 trading days, and 65 trading days. The base-5 geometric spacing provides convenient multi-scale coverage at 30-minute resolution, analogous to the daily/weekly/monthly structure in \citet{Corsi2009} but extended to higher frequencies. For each lag $k \in \mathcal{L}$, the HAR feature is the trailing rolling mean of the preprocessed target:
\begin{equation}
    \overline{RV}_t^{(k)} = \frac{1}{k}\sum_{i=0}^{k-1} \widetilde{RV}_{t-1-i},
    \label{eq:har_ma}
\end{equation}
where the shift by one ensures no contemporaneous information enters the feature. The standard HAR regression model takes the form
\begin{equation}
    \widetilde{RV}_{t+1} = \beta_0 + \sum_{k \in \mathcal{L}} \beta_k \, \overline{RV}_t^{(k)} + \varepsilon_{t+1}.
    \label{eq:har_model}
\end{equation}

When exogenous variables are included, the same HAR lag structure is applied to each additional preprocessed column, generating $|\mathcal{L}| = 6$ rolling-mean features per variable.

\subsubsection{PCA Features}
\label{sec:pca_features}

As an alternative to the fixed HAR lag structure, we consider a data-driven approach that serves two purposes: (i) it lets the data determine the optimal lag weighting rather than imposing the HAR rolling-mean structure, and (ii) it reduces dimensionality when many exogenous variables are included. We first construct approximately 20 log-spaced raw lags spanning the range $[1, 3125]$, providing dense coverage at short horizons and sparser coverage at longer horizons. For each variable, individual shifted lags $x_{t-\ell}$ for $\ell \in \mathcal{L}_{\mathrm{pca}}$ are computed (rather than rolling means as in the HAR approach).

These raw lag features are then projected onto a low-dimensional subspace via Principal Component Analysis (PCA). Specifically, on each rolling training window of length $T_{\mathrm{train}}$, we fit PCA with $n_c = 5$ components using a randomized SVD solver:
\begin{equation}
    \mathbf{z}_t = \mathbf{W}^{\top}(\mathbf{x}_t - \bar{\mathbf{x}}),
    \label{eq:pca}
\end{equation}
where $\mathbf{W} \in \mathbb{R}^{p \times n_c}$ is the projection matrix estimated from the training buffer and $\bar{\mathbf{x}}$ is the training-set mean. The PCA transform is re-estimated each time the model is refit, ensuring the projection adapts to evolving autocorrelation structures.

\subsubsection{Exogenous Variable Subgroups}
\label{sec:subgroups}

To assess the marginal contribution of different information sources, we organize exogenous variables into thematic subgroups. Each subgroup augments the target-only baseline (HAR or PCA features of $\widetilde{RV}$ alone) with additional features derived from the subgroup's variables using the same lag structure. The subgroups are:

\begin{itemize}
    \item \textbf{Baseline}: HAR (or PCA) features of $\widetilde{RV}$ only --- no exogenous variables.
    \item \textbf{Moments}: Higher-order return moments and related measures (bipower variation, realized semi-variance, autocovariance, absolute returns, cubic and quartic returns).
    \item \textbf{Liquidity}: Order flow and market microstructure variables (volume, turnover, buy/sell turnover, effective spread, quoted spread, number of observations).
    \item \textbf{Market (EW)}: Equal-weighted cross-sectional stock factor aggregates (returns, bipower, semi-variance, absolute returns).
    \item \textbf{Market (VW)}: Value-weighted analogues of the equal-weighted factors.
    \item \textbf{Sentiment}: Social media sentiment metrics (StockTwits attention, sentiment score, sentiment count).
    \item \textbf{Implied Volatility}: Options-implied measures (VIX, VVIX, VIX3M).
    \item \textbf{Volatility Demand}: Options order flow measures (SPX open/close, all options open/close, open only).
    \item \textbf{All Features}: The union of all exogenous variables above.
\end{itemize}

\subsection{Models}
\label{sec:models}

We consider five models spanning parametric linear, ensemble tree-based, and naive approaches. All models are trained on the preprocessed (diurnally adjusted, transformed, winsorized) scale. Key differences in refit frequency reflect computational cost: ridge regression admits a closed-form solution and is refit at every time step, while tree ensembles are more expensive and are refit every 5 steps. Section~\ref{sec:composed} additionally describes a family of \emph{composed} models --- a baseline whose residuals are modeled by a second regressor, optionally through a learned feature map --- of which the five models here are the degenerate single-stage case.

\subsubsection{Naive Baseline}
\label{sec:naive}

The naive baseline predicts $\widetilde{RV}_{t+1}$ using the HAR rolling mean at the 125-period horizon (approximately 2.6 trading days):
\begin{equation}
    \hat{y}_{t+1}^{\mathrm{naive}} = \overline{RV}_t^{(125)}.
    \label{eq:naive}
\end{equation}
This provides a non-trivial benchmark that reflects the strong persistence of realized volatility.

\subsubsection{Ridge Regression}
\label{sec:ridge}

Ridge regression augments ordinary least squares with an $\ell_2$ penalty on the coefficient vector:
\begin{equation}
    \hat{\boldsymbol{\beta}} = \arg\min_{\boldsymbol{\beta}} \left\{ \|\mathbf{y} - \mathbf{X}\boldsymbol{\beta}\|_2^2 + \alpha \|\boldsymbol{\beta}\|_2^2 \right\},
    \label{eq:ridge}
\end{equation}
with regularization parameter $\alpha = 1.0$. Input features are standardized using a rolling robust scaler based on the median and interquartile range (IQR) estimated from the current training buffer. Ridge is refit at every time step via the closed-form solution.

\subsubsection{XGBoost}
\label{sec:xgboost}

We employ XGBoost \citep{Chen2016} as a gradient-boosted tree ensemble with the histogram-based splitting method. Two hyperparameter configurations are evaluated: a default configuration (maximum depth 3, 100 boosting rounds, learning rate 0.1), and a tuned configuration obtained via the Optuna procedure described in Section~\ref{sec:tuning}. The model is refit every 5 time steps.

\subsubsection{LightGBM}
\label{sec:lightgbm}

LightGBM \citep{Ke2017} is similarly evaluated in default (matching XGBoost defaults) and tuned configurations. The model is refit every 5 time steps.

\subsubsection{Random Forest}
\label{sec:rf}

A random forest ensemble is evaluated in default (100 trees, maximum depth 3) and tuned configurations. Like the other tree-based models, it is refit every 5 time steps.

\subsubsection{Hyperparameter Tuning}
\label{sec:tuning}

For each tree-based model we perform Bayesian hyperparameter optimization using Optuna \citep{Akiba2019} with the Tree-structured Parzen Estimator (TPE) sampler. The objective is to minimize QLIKE (Section~\ref{sec:qlike}) on the full walk-forward backtest. To support distributed parallel trial evaluation across HPC clusters, we enable the \texttt{constant\_liar} option so that multiple workers can draw independent candidates from the same study without collision. No pruner is configured; every trial runs to completion.

\paragraph{Search spaces.}
The per-model search spaces are summarized in Table~\ref{tab:tune_search_space}. Ranges are chosen to bracket commonly-reported values in the gradient-boosting literature while allowing substantial deviation from the defaults. Integer dimensions are sampled uniformly; learning rates and regularization strengths are sampled on a log scale to cover several orders of magnitude.

\begin{table}[H]
\centering
\caption{Hyperparameter search spaces per model. ``Log'' indicates log-uniform sampling; ``step'' denotes the integer step size.}
\label{tab:tune_search_space}
\begin{tabular}{lllc}
\toprule
Model & Parameter & Range & Scale \\
\midrule
\multirow{6}{*}{RF}
  & n\_estimators      & $[100,\,1000]$, step 100 & linear \\
  & max\_depth         & $[3,\,20]$               & linear \\
  & min\_samples\_leaf & $[1,\,50]$               & log    \\
  & min\_samples\_split& $[2,\,20]$               & linear \\
  & max\_features      & $[0.3,\,1.0]$            & linear \\
  & max\_samples       & $[0.5,\,1.0]$            & linear \\
\midrule
\multirow{8}{*}{XGB}
  & n\_estimators      & $[100,\,1000]$, step 100 & linear \\
  & max\_depth         & $[3,\,12]$               & linear \\
  & learning\_rate     & $[0.01,\,0.3]$           & log    \\
  & min\_child\_weight & $[1,\,20]$               & linear \\
  & subsample          & $[0.5,\,1.0]$            & linear \\
  & colsample\_bytree  & $[0.3,\,1.0]$            & linear \\
  & reg\_alpha         & $[10^{-8},\,10]$         & log    \\
  & reg\_lambda        & $[10^{-8},\,10]$         & log    \\
\midrule
\multirow{9}{*}{LGBM}
  & n\_estimators      & $[100,\,1000]$, step 100 & linear \\
  & max\_depth         & $[3,\,12]$               & linear \\
  & learning\_rate     & $[0.01,\,0.3]$           & log    \\
  & num\_leaves        & $[15,\,127]$             & linear \\
  & min\_child\_samples& $[5,\,100]$              & linear \\
  & subsample          & $[0.5,\,1.0]$            & linear \\
  & colsample\_bytree  & $[0.3,\,1.0]$            & linear \\
  & reg\_alpha         & $[10^{-8},\,10]$         & log    \\
  & reg\_lambda        & $[10^{-8},\,10]$         & log    \\
\bottomrule
\end{tabular}
\end{table}

\paragraph{Trial evaluation.}
Each trial is a full walk-forward backtest over the out-of-sample window, partitioned into $C=100$ contiguous chunks executed in parallel. Each chunk receives a 24{,}000-observation prefix of training data (approximately 500 trading days at 48 bars per day) and produces predictions on its assigned OOS slice. Chunk-level results are concatenated and scored via QLIKE on the raw-variance scale after inverting the preprocessing pipeline. The resulting scalar is reported back to the Optuna study and used by the TPE sampler to propose subsequent candidates. Trials are distributed across two HPC clusters (UCLA Hoffman2 and USC CARC Discovery) with SGE and SLURM array jobs respectively; the shared Optuna storage ensures that candidate suggestions reflect all completed trials regardless of where they ran.

\medskip

\noindent\textit{Categorical time features.} All three tree-based models (XGBoost, LightGBM, random forest) additionally receive day-of-week and hour-of-day as categorical feature columns appended to the HAR or PCA features. Ridge regression does not receive these features, as the diurnal adjustment (Stage 1) already removes time-of-day effects from the input series.

\subsection{Composed Models: Residualization, Manifold Features, and Local Regression}
\label{sec:composed}

The five models of Section~\ref{sec:models} are \emph{single-stage}: a regressor maps the feature vector $\mathbf{x}_t$ to a forecast of $y_t \equiv \widetilde{RV}_{t}$ directly. We additionally consider \emph{composed} models built from a common three-stage template, which nests the single-stage case as a degenerate configuration. Writing the forecast as

\begin{equation}
    \hat{y}_t \;=\; \underbrace{g_t(\mathbf{x}_t)}_{\text{baseline}} \;+\; \underbrace{h_t\!\left(\Phi_t(\mathbf{v}_t)\right)}_{\text{residual correction}},
    \label{eq:multistage}
\end{equation}

the three stages are:

\begin{enumerate}
    \item \textbf{Residualizer} $g_t$ --- a baseline regressor fit on the current training window, whose prediction is subtracted from the target. The quantity passed downstream is the residual series $e_s = y_s - g_t(\mathbf{x}_s)$. Setting $g_t \equiv 0$ (so that $e_s = y_s$) recovers the single-stage case.
    \item \textbf{Feature map} $\Phi_t$ --- an optional unsupervised transform, re-estimated on a slower cadence, mapping a representation $\mathbf{v}_t$ of recent residuals into a learned coordinate system. When omitted, the downstream regressor consumes $\mathbf{x}_t$ directly.
    \item \textbf{Regressor} $h_t$ --- a supervised model fit to predict the residual from the (possibly transformed) features.
\end{enumerate}

Three configurations are of interest, summarized in Table~\ref{tab:composed_configs}. All share the identical data-preparation path of Sections~\ref{sec:preprocessing}--\ref{sec:features} and the identical walk-forward protocol of Section~\ref{sec:backtest}; they differ only in which stages are active. Strict causality is preserved in every configuration: at step $t$ only $\{\mathbf{x}_s, y_s\}_{s<t}$ and the contemporaneous feature vector $\mathbf{x}_t$ are read.

\begin{table}[H]
\centering
\caption{Stage configurations of the composed-model template, Equation~\eqref{eq:multistage}. ``---'' denotes an inactive (identity or omitted) stage.}
\label{tab:composed_configs}
\begin{tabular}{llll}
\toprule
Configuration & Residualizer $g_t$ & Feature map $\Phi_t$ & Regressor $h_t$ \\
\midrule
Single-stage      & ---            & ---                & Ridge / XGBoost / kNN / \ldots \\
Stacking          & Ridge          & ---                & kNN \\
Spectral kNN      & Ridge          & Laplacian eigenmap & kNN \\
\bottomrule
\end{tabular}
\end{table}

\subsubsection{kNN Regression with a Self-Tuning Bandwidth}
\label{sec:knn}

The regressor stage in both composed configurations is $k$-nearest-neighbour regression with a Gaussian kernel. Given a query point $\mathbf{q}$, a reference set $\{(\boldsymbol{\phi}_i, z_i)\}$, and the index set $\mathcal{N}_k(\mathbf{q})$ of its $k$ nearest reference points under the Euclidean metric, the prediction is the kernel-weighted average

\begin{equation}
    h(\mathbf{q}) \;=\; \frac{\sum_{i \in \mathcal{N}_k(\mathbf{q})} w_i \, z_i}{\sum_{i \in \mathcal{N}_k(\mathbf{q})} w_i},
    \qquad
    w_i \;=\; \exp\!\left(-\frac{\|\mathbf{q} - \boldsymbol{\phi}_i\|_2^2}{2\,\sigma_{\mathbf{q}}^2}\right).
    \label{eq:knn_gaussian}
\end{equation}

The bandwidth is \emph{self-tuning}: rather than fixing a global $\sigma$, we set it per query to the median distance to that query's own neighbours,

\begin{equation}
    \sigma_{\mathbf{q}} \;=\; \mathrm{median}_{\,i \in \mathcal{N}_k(\mathbf{q})} \, \|\mathbf{q} - \boldsymbol{\phi}_i\|_2,
    \label{eq:knn_bandwidth}
\end{equation}

in the spirit of the local-scaling rule of \citet{ZelnikManor2004}. This makes the weighting profile invariant to the local density of the reference set: the median-distance neighbour always receives weight $e^{-1/2} \approx 0.607$ relative to a coincident point, whether the query lands in a densely populated calm-market region or in a sparse crisis region. A fixed bandwidth would instead collapse to near-uniform weighting in dense regions and to nearest-neighbour weighting in sparse ones. Because Equation~\eqref{eq:knn_gaussian} normalizes by $\sum_i w_i$, any common scale factor on the weights cancels, so only the \emph{shape} induced by $\sigma_{\mathbf{q}}$ matters.

We set $k = 25$ throughout. As a standalone (single-stage) model, kNN regresses $y_t$ on the rolling-robust-scaled HAR feature matrix; because the reference set is refit merely by rebuilding a spatial index, it is refit at every step.

\subsubsection{Spectral kNN}
\label{sec:spectral_knn}

The spectral-kNN model is motivated by the hypothesis that what a linear HAR specification fails to explain is not white noise but a low-dimensional, recurring \emph{regime} structure --- that residual trajectories through a volatility crisis resemble one another more than they resemble residual trajectories through a calm period, in a way no fixed lag basis captures. The model therefore (i) strips out the linearly predictable component with a ridge baseline, (ii) represents recent history as a delay-embedded trajectory of residuals, (iii) learns an unsupervised low-dimensional coordinate system on those trajectories via a graph Laplacian, and (iv) forecasts the next residual by local averaging among trajectories that are neighbours \emph{in that learned coordinate system}. Steps (ii)--(iii) are the manifold-learning analogue of the PCA feature construction of Section~\ref{sec:pca_features}, with a nonlinear, locality-preserving map replacing the linear projection.

\paragraph{Stage 1: Ridge residualization.}
At each test step $t$, a ridge regression with $\alpha = 1.0$ (Section~\ref{sec:ridge}) is refit on the trailing training window $\mathcal{T}_t = \{t - T_{\mathrm{train}}, \ldots, t-1\}$, yielding coefficients $\hat{\boldsymbol{\beta}}^{(t)}$. The residual at any index $s$ is

\begin{equation}
    e^{(t)}_s \;=\; y_s - \mathbf{x}_s^{\top}\hat{\boldsymbol{\beta}}^{(t)}.
    \label{eq:residual}
\end{equation}

Refitting every step is affordable because ridge admits a closed-form solution. The baseline forecast entering Equation~\eqref{eq:multistage} is $g_t(\mathbf{x}_t) = \mathbf{x}_t^{\top}\hat{\boldsymbol{\beta}}^{(t)}$.

\paragraph{Stage 2: Delay-embedded residual views.}
The remaining stages are re-estimated only every $F = 240$ steps (five trading days) to amortize their cost. At a refit step $t_r$, we form the training residual series $\{e^{(t_r)}_s\}_{s \in \mathcal{T}_{t_r}}$ and slide a window of length $W = 960$ bars (twenty trading days) across it. Re-indexing the window locally as $u = 1, \ldots, T_{\mathrm{train}}$, each position $j \in \{W+1, \ldots, T_{\mathrm{train}}\}$ yields a \emph{view} and a scalar target,

\begin{equation}
    \mathbf{v}_j \;=\; \bigl(e_{j-W},\, e_{j-W+1},\, \ldots,\, e_{j-1}\bigr) \in \mathbb{R}^{W},
    \qquad
    z_j \;=\; e_j,
    \label{eq:views}
\end{equation}

stacked into a view matrix $V \in \mathbb{R}^{N \times W}$ with $N = T_{\mathrm{train}} - W$ rows. This is a delay-coordinate (Takens-style) reconstruction of the residual process \citep{Takens1981}: a point of $V$ is an entire twenty-day residual \emph{trajectory}, not a single observation, and the pairing $(\mathbf{v}_j, z_j)$ is by construction a one-step-ahead prediction problem. With the production setting $T_{\mathrm{train}} = 500 \times 48 = 24{,}000$ bars, this gives $N = 23{,}040$ trajectories in $\mathbb{R}^{960}$.

\paragraph{Stage 3: Laplacian eigenmap.}
We embed the view matrix with a Laplacian eigenmap \citep{Belkin2003,vonLuxburg2007}. First, a symmetric $k$-nearest-neighbour \emph{connectivity} graph is built over the $N$ views with $k_{\mathrm{graph}} = 10$: writing $C_{ij} = 1$ if $\mathbf{v}_j$ is among the $k_{\mathrm{graph}}$ nearest neighbours of $\mathbf{v}_i$ (and $C_{ii}=1$), the affinity is the symmetrization

\begin{equation}
    A \;=\; \tfrac{1}{2}\left(C + C^{\top}\right),
    \label{eq:affinity}
\end{equation}

so mutual neighbours carry weight $1$ and one-directional neighbours weight $\tfrac{1}{2}$. Note that the affinity is \emph{binary} (up to the symmetrization) rather than a Gaussian kernel of the distances: the graph encodes only \emph{who is near whom}, discarding the magnitude of the distance. This is deliberate --- in $\mathbb{R}^{960}$ raw Euclidean distances concentrate and carry little information, whereas the rank ordering that defines the neighbour sets remains meaningful.

Given the degree matrix $D = \mathrm{diag}(A\mathbf{1})$, the symmetric normalized Laplacian is

\begin{equation}
    L_{\mathrm{sym}} \;=\; I - D^{-1/2} A D^{-1/2},
    \label{eq:laplacian}
\end{equation}

and the embedding coordinates are the eigenvectors associated with its smallest eigenvalues, rescaled by $D^{-1/2}$ to recover the eigenvectors of the random-walk Laplacian $L_{\mathrm{rw}} = I - D^{-1}A$:

\begin{equation}
    L_{\mathrm{sym}}\,\mathbf{u}_m \;=\; \lambda_m\,\mathbf{u}_m,
    \qquad 0 = \lambda_0 \leq \lambda_1 \leq \cdots,
    \qquad
    \boldsymbol{\psi}_m \;=\; D^{-1/2}\mathbf{u}_m .
    \label{eq:eigenproblem}
\end{equation}

Retaining $d = 8$ coordinates gives the training embedding $\Psi \in \mathbb{R}^{N \times d}$, whose $j$-th row $\boldsymbol{\phi}_j = (\psi_1(j), \ldots, \psi_d(j))$ is the spectral representation of trajectory $\mathbf{v}_j$. The eigenproblem is solved with the ARPACK implicitly-restarted Lanczos solver in shift-invert mode \citep{Lehoucq1996}, which is well suited to extracting a few extreme eigenpairs of a large sparse matrix.

The variational characterization makes the objective explicit: the eigenvectors minimize $\sum_{i,j} A_{ij}\,(\psi(i) - \psi(j))^2$ subject to normalization and orthogonality constraints, so the embedding places trajectories that are graph neighbours close together while spreading apart trajectories connected only through long paths. It is a \emph{locality-preserving} map, and --- unlike PCA --- makes no claim to preserve global geometry. It is also entirely \emph{unsupervised}: targets $z_j$ never enter Stage 3, so the coordinate system is a description of the residual dynamics, not a fit to them.

\paragraph{Stage 4: Nystr\"om out-of-sample extension.}
Equation~\eqref{eq:eigenproblem} defines coordinates only for the $N$ training views; it provides no formula for a new trajectory. We extend the map out of sample with the weighted-kNN Nystr\"om approximation \citep{Bengio2004,Williams2001}. For a test view $\mathbf{v}_\ast$, let $\mathcal{N}_{k_{\mathrm{graph}}}(\mathbf{v}_\ast)$ be its $k_{\mathrm{graph}}$ nearest training views with distances $\delta_i = \|\mathbf{v}_\ast - \mathbf{v}_i\|_2$. Then

\begin{equation}
    \Phi(\mathbf{v}_\ast) \;=\; \frac{\sum_{i \in \mathcal{N}_{k_{\mathrm{graph}}}} \tilde{w}_i \, \boldsymbol{\phi}_i}{\sum_{i \in \mathcal{N}_{k_{\mathrm{graph}}}} \tilde{w}_i},
    \qquad
    \tilde{w}_i = \exp\!\left(-\frac{\delta_i^2}{2\tilde{\sigma}_\ast^2}\right),
    \qquad
    \tilde{\sigma}_\ast = \mathrm{median}_{\,i}\,\delta_i,
    \label{eq:nystrom}
\end{equation}

i.e.\ the new point inherits a convex combination of its neighbours' spectral coordinates, under the same self-tuning Gaussian kernel used by the regressor (Equations~\eqref{eq:knn_gaussian}--\eqref{eq:knn_bandwidth}). If all weights underflow to zero, the extension degrades gracefully to the unweighted neighbour mean. This step is what makes the embedding usable in a forecasting loop at all: the training-side eigendecomposition is computed once per refit, and each of the $F$ subsequent test points is embedded at the cost of a single $k$-nearest-neighbour query.

\paragraph{Stage 5: Regression in embedding space and recombination.}
A kNN regressor with $k = 25$ and the self-tuning Gaussian weights of Section~\ref{sec:knn} is fit on the pairs $\{(\boldsymbol{\phi}_j, z_j)\}_{j=1}^{N}$ --- that is, in the $d = 8$ dimensional spectral space, not in the $960$-dimensional view space. At test step $t$ the view $\mathbf{v}_t = (e^{(t)}_{t-W}, \ldots, e^{(t)}_{t-1})$ is formed from residuals against the \emph{current} step's ridge fit, embedded via Equation~\eqref{eq:nystrom}, and passed to the regressor; its output is the residual correction added to the ridge baseline in Equation~\eqref{eq:multistage}. Algorithm~\ref{alg:spectral_knn} states the complete procedure.

\begin{algorithm}[H]
\caption{Spectral kNN --- one walk-forward pass}
\label{alg:spectral_knn}
\begin{algorithmic}[1]
\Require features $X$, target $y$, train window $T_{\mathrm{train}}$, ridge penalty $\alpha$, view length $W$, embedding dimension $d$, graph degree $k_{\mathrm{graph}}$, regression degree $k$, refit cadence $F$
\For{$t = T_{\mathrm{train}}, \ldots, n-1$}
    \State $i \gets t - T_{\mathrm{train}}$ \Comment{test-step counter}
    \State $\hat{\boldsymbol{\beta}} \gets \textsc{RidgeFit}\bigl(X[t-T_{\mathrm{train}}\!:\!t],\; y[t-T_{\mathrm{train}}\!:\!t];\; \alpha\bigr)$ \Comment{Stage 1, every step}
    \If{$i \bmod F = 0$} \Comment{Stages 2--3, every $F$ steps}
        \State $e_s \gets y_s - \mathbf{x}_s^{\top}\hat{\boldsymbol{\beta}}$ for $s \in [t-T_{\mathrm{train}}, t)$
        \If{$T_{\mathrm{train}} < W + 1$} \Comment{window too short to form a view}
            \State $\Psi, \textsc{knn} \gets \varnothing$ \Comment{fall back to the baseline alone}
        \Else
            \State $V \gets \bigl[\,\mathbf{v}_j = (e_{j-W}, \ldots, e_{j-1})\,\bigr]_{j=W+1}^{T_{\mathrm{train}}}$, \; $z_j \gets e_j$
            \State $A \gets \tfrac{1}{2}(C + C^{\top})$ from the $k_{\mathrm{graph}}$-NN connectivity graph of $V$
            \State $\Psi \gets$ bottom-$d$ eigenvectors of $L_{\mathrm{sym}} = I - D^{-1/2}AD^{-1/2}$, rescaled by $D^{-1/2}$
            \State cache a $k_{\mathrm{graph}}$-NN index over $V$ for Nystr\"om extension
            \State $\textsc{knn} \gets \textsc{KNNFit}(\Psi, z;\; k,\; \text{Gaussian self-tuning weights})$
        \EndIf
    \EndIf
    \State $\hat{y}^{\mathrm{base}} \gets \mathbf{x}_t^{\top}\hat{\boldsymbol{\beta}}$
    \If{$\textsc{knn} = \varnothing$}
        \State $\hat{y}_t \gets \hat{y}^{\mathrm{base}}$ \Comment{degenerate case: ridge forecast, uncorrected}
    \Else
        \State $\mathbf{v}_\ast \gets \bigl(y_s - \mathbf{x}_s^{\top}\hat{\boldsymbol{\beta}}\bigr)_{s=t-W}^{t-1}$ \Comment{test view, current-step baseline}
        \State $\boldsymbol{\phi}_\ast \gets \textsc{Nystr\"om}(\mathbf{v}_\ast)$ \Comment{Equation~\eqref{eq:nystrom}}
        \State $\hat{y}_t \gets \hat{y}^{\mathrm{base}} + \textsc{knn}.\textsc{Predict}(\boldsymbol{\phi}_\ast)$
    \EndIf
\EndFor
\end{algorithmic}
\end{algorithm}

\paragraph{Hyperparameters.}
Table~\ref{tab:spectral_knn_params} lists the production configuration. Unlike the tree models, spectral kNN is not tuned by Optuna; the values are set from the timescale interpretation of each parameter rather than by search.

\begin{table}[H]
\centering
\caption{Spectral kNN hyperparameters. Bar counts convert at 48 bars per 24-hour trading day (Section~\ref{sec:rv_definition}).}
\label{tab:spectral_knn_params}
\begin{tabular}{llll}
\toprule
Symbol & Parameter & Value & Interpretation \\
\midrule
$\alpha$            & \texttt{ridge\_alpha}      & $1.0$   & Baseline $\ell_2$ penalty; matches Section~\ref{sec:ridge} \\
$T_{\mathrm{train}}$& \texttt{train\_window}     & $500$ d $= 24{,}000$ bars & Rolling training window \\
$W$                 & \texttt{view\_window}      & $960$ bars ($20$ d) & Length of a residual trajectory \\
$d$                 & \texttt{embedding\_dim}    & $8$     & Retained Laplacian coordinates \\
$k_{\mathrm{graph}}$& \texttt{graph\_k}          & $10$    & Affinity-graph degree; also the Nystr\"om degree \\
$k$                 & \texttt{neighbor\_k}       & $25$    & Neighbours for the embedding-space regression \\
$F$                 & \texttt{refit\_frequency}  & $240$ steps ($5$ d) & Embedding + regressor refit cadence \\
\bottomrule
\end{tabular}
\end{table}

\paragraph{Computational cost.}
The dominant term is Stage 3. Each refit builds a $k_{\mathrm{graph}}$-NN graph over $N = 23{,}040$ points in $\mathbb{R}^{960}$ and extracts eight eigenpairs of the resulting sparse $N \times N$ Laplacian; the graph construction is $O(N^2 W)$ in the worst case and the Lanczos iteration is $O(d \cdot \mathrm{nnz}(A))$ per pass, with $\mathrm{nnz}(A) = O(N k_{\mathrm{graph}})$. The refit cadence $F = 240$ amortizes this over five trading days of forecasts, and the per-step marginal cost is only a ridge solve plus two $k$-NN queries. This is why $F$ is two orders of magnitude larger here than for the tree models, which refit every five steps.

\paragraph{Design notes and caveats.}
Several properties of the construction are worth making explicit.

\begin{itemize}
    \item \textit{Baseline vintage is not synchronized across stages.} The training views at refit step $t_r$ are residuals against $\hat{\boldsymbol{\beta}}^{(t_r)}$, whereas the test view at step $t > t_r$ is built from residuals against the fresher $\hat{\boldsymbol{\beta}}^{(t)}$. Between refits the two residual definitions therefore drift apart by $\mathbf{x}_s^{\top}(\hat{\boldsymbol{\beta}}^{(t)} - \hat{\boldsymbol{\beta}}^{(t_r)})$. The alternative --- freezing the baseline at $t_r$ for the test view --- would keep the residual definition internally consistent but would leave the \emph{baseline forecast} itself stale for up to $F$ steps. The present choice prioritizes baseline freshness; the sensitivity of forecasts to this choice has not been measured, and doing so is left for future work.
    \item \textit{The trivial eigenvector.} The eigenvector associated with $\lambda_0 = 0$ is constant on each connected component of the graph and carries no information about within-component geometry. Implementations differ in whether it is discarded; if it is retained, one of the $d$ coordinates is uninformative and the effective embedding dimension is $d-1$. The concrete count should be verified against the pinned \texttt{scikit-learn} version before the coordinates are interpreted individually.
    \item \textit{Graph connectivity.} With $k_{\mathrm{graph}} = 10$ over $23{,}040$ points there is no guarantee that the affinity graph is connected. A disconnected graph has one zero eigenvalue per component, so the leading coordinates become component-indicator functions rather than smooth geometric coordinates --- a coarse regime partition rather than a manifold parameterization. Whether this occurs in practice is an empirical question about the data, and the number of near-zero eigenvalues per refit is the natural diagnostic.
    \item \textit{Overlapping views induce strong dependence.} Consecutive views share $W-1$ of their $W$ coordinates, so neighbouring rows of $V$ are nearly identical by construction and the graph is guaranteed to link each view to its temporal neighbours. Some of the recovered structure is therefore an artifact of the sliding window rather than evidence of recurring regimes; a meaningful embedding is one in which \emph{temporally distant} views are also linked.
    \item \textit{Distance concentration.} Euclidean distances in $\mathbb{R}^{960}$ are poorly discriminating in the absolute sense. The design mitigates but does not eliminate this: the affinity graph uses only neighbour \emph{ranks}, and the final regression runs in $\mathbb{R}^{8}$ where distances are again well behaved. Only the neighbour queries --- graph construction and Nystr\"om extension --- operate in the high-dimensional space.
    \item \textit{Feature scaling.} As with plain kNN, the feature matrix is rolling-robust-scaled whole-series before the backtest (\texttt{prescale}), so the ridge residuals --- and hence the views --- are on a comparable per-dimension scale and the Euclidean metric is not dominated by an arbitrarily-scaled column.
    \item \textit{Graceful degradation.} If the training window is too short to form even one view ($T_{\mathrm{train}} < W + 1$), the feature and regressor stages are skipped and the model emits the ridge baseline alone. Spectral kNN thus never performs worse than its baseline by construction failure --- only by a genuinely harmful residual correction.
\end{itemize}

\subsection{Walk-Forward Backtesting}
\label{sec:backtest}

All models are evaluated using a strictly out-of-sample walk-forward procedure. At each test period $t$, the model is trained (or updated) on a rolling window of $T_{\mathrm{train}} = 500$ consecutive observations immediately preceding $t$. With 48 bars per 24-hour trading day, this corresponds to approximately 10.4 trading days of training data. The procedure is:

\begin{enumerate}
    \item \textbf{Initialize}: Fit the model on observations $\{t - T_{\mathrm{train}}, \ldots, t-1\}$.
    \item \textbf{Predict}: Generate forecast $\hat{y}_{t}$ using features $\mathbf{x}_t$.
    \item \textbf{Observe}: Record the realized target $y_t$.
    \item \textbf{Update}: Add $(\mathbf{x}_t, y_t)$ to the rolling buffer, drop the oldest observation, and refit the model at the model-specific refit frequency (every step for ridge and kNN, every 5 steps for tree models, every 240 steps for the spectral embedding and its downstream regressor --- see Table~\ref{tab:spectral_knn_params}; the ridge baseline within spectral kNN still refits every step).
    \item \textbf{Advance}: Increment $t$ and return to step 2.
\end{enumerate}

For ridge regression, feature scaling parameters (median and IQR) are updated at every step using the rolling robust scaler fit on the current training buffer. For PCA-based configurations, the projection matrix $\mathbf{W}$ is re-estimated at each refit. This ensures all model components adapt to non-stationarities while maintaining strict temporal separation between training and test data.

The backtest is parallelized across 100 chunks via SLURM array jobs, with each chunk processing a contiguous subset of test indices.

\subsection{Evaluation}
\label{sec:evaluation}

\subsubsection{Inverse Transform and Duan's Smearing Estimator}
\label{sec:duan}

Models are trained on the preprocessed scale (after diurnal adjustment, semantic transform, and winsorization). To evaluate forecasts against raw realized volatility, predictions must be mapped back through the inverse of the preprocessing pipeline. Of the three stages, only Stages 1 (diurnal adjustment) and 2 (semantic transform) are invertible; Stage 3 (winsorization) is lossy and is not reversed.

For the primary target $RV$, the semantic transform is the square root ($\sqrt{\cdot}$), so the na\"ive inverse would simply square the forecast: $\hat{y}_t^2$. However, because $\mathbb{E}[\sqrt{X}]^2 \neq \mathbb{E}[X]$ by Jensen's inequality, squaring a forecast of $\sqrt{RV}$ produces a downward-biased estimate of $RV$. Duan's smearing estimator \citep{Duan1983} corrects this bias by adding the in-sample residual variance before inverting:
\begin{equation}
    \widehat{RV}_t^{\mathrm{raw}} = \left(\hat{y}_t^2 + \hat{\sigma}_{\varepsilon}^2\right) \cdot \bar{RV}_{t \mid s}^{\mathrm{baseline}},
    \label{eq:duan}
\end{equation}
where $\hat{\sigma}_{\varepsilon}^2 = \frac{1}{N}\sum_{i}(y_i - \hat{y}_i)^2$ is the in-sample mean squared residual (estimating the variance lost by working on the transformed scale) and $\bar{RV}_{t \mid s}^{\mathrm{baseline}}$ is the diurnal baseline factor at time $t$ that reverses Stage 1.

\subsubsection{QLIKE Loss}
\label{sec:qlike}

Our primary evaluation metric is the quasi-likelihood (QLIKE) loss function, which is the robust loss function for volatility forecast comparison under general loss functions \citep{Patton2011}:
\begin{equation}
    L_{\mathrm{QLIKE}} = \frac{1}{N}\sum_{t=1}^{N} \left[ \frac{RV_t^{\mathrm{raw}}}{\widehat{RV}_t^{\mathrm{raw}}} - \log\!\left(\frac{RV_t^{\mathrm{raw}}}{\widehat{RV}_t^{\mathrm{raw}}}\right) - 1 \right].
    \label{eq:qlike}
\end{equation}
QLIKE is strictly consistent for the conditional mean under the assumption that realized volatility is an unbiased estimator of integrated variance \citep{Patton2011}. It penalizes under-prediction more heavily than over-prediction, reflecting the asymmetric costs relevant in risk management applications.

\subsubsection{Additional Metrics}
\label{sec:additional_metrics}

We also report mean squared error (MSE) and mean absolute error (MAE) computed on the preprocessed (adjusted) scale. To limit the influence of extreme outliers, winsorized variants (clipping errors to the 5th--95th percentile range) are computed for all metrics. Out-of-sample $R^2$ is defined relative to the naive baseline:
\begin{equation}
    R^2_{\mathrm{OOS}} = 1 - \frac{\mathrm{MSE}_{\mathrm{model}}}{\mathrm{MSE}_{\mathrm{naive}}},
    \label{eq:oos_r2}
\end{equation}
where positive values indicate improvement over the naive forecast.

\subsection{Experimental Design}
\label{sec:experimental_design}

We report results from three complementary analyses, each designed to isolate a specific dimension of the forecasting problem:

\begin{enumerate}
    \item \textbf{Ridge--HAR Subgroup Analysis.} Ridge regression with HAR features is estimated for each exogenous subgroup (Section~\ref{sec:subgroups}). By holding the model fixed (ridge) and varying the information set, this analysis isolates the marginal predictive content of each information source beyond the target's own history.

    \item \textbf{HAR Model Comparison.} Using only the target HAR features (baseline subgroup, no exogenous variables), we compare all five models --- naive, ridge, XGBoost, LightGBM, and random forest. By holding the features fixed and varying the model, this analysis assesses whether nonlinear models capture dynamics beyond the linear HAR specification.

    \item \textbf{Ridge--PCA Subgroup Analysis.} Ridge regression with PCA features is estimated for each subgroup, paralleling the HAR analysis. By replacing the hand-crafted HAR lag structure with data-driven PCA compression while keeping the model fixed, this analysis evaluates whether flexible lag weighting and dimensionality reduction offer complementary predictive power.
\end{enumerate}
